% This is samplepaper.tex, a sample chapter demonstrating the
% LLNCS macro package for Springer Computer Science proceedings;
% Version 2.21 of 2022/01/12
%
\documentclass[runningheads]{llncs}
%
\usepackage[T1]{fontenc}
% T1 fonts will be used to generate the final print and online PDFs,
% so please use T1 fonts in your manuscript whenever possible.
% Other font encondings may result in incorrect characters.
%
\usepackage{graphicx}
% Used for displaying a sample figure. If possible, figure files should
% be included in EPS format.
%
% If you use the hyperref package, please uncomment the following two lines
% to display URLs in blue roman font according to Springer's eBook style:
%\usepackage{color}
%\renewcommand\UrlFont{\color{blue}\rmfamily}
%\urlstyle{rm}
%
\begin{document}
%
\title{Individual Report}
%
%\titlerunning{Abbreviated paper title}
% If the paper title is too long for the running head, you can set
% an abbreviated paper title here
%
\author{TRailMix: Lena Friedek\inst{2}}
%
%\authorrunning{F. Author et al.}
% First names are abbreviated in the running head.
% If there are more than two authors, 'et al.' is used.
%
\institute{
Universität Potsdam,
\email{lena.friedek@uni-potsdam.de}}
%
\maketitle              % typeset the header of the contribution
%
%\begin{abstract}
%The abstract should briefly summarize the contents of the paper in
%150--250 words.
%
%\keywords{clingo  \and flatland \and MAPF \and Rescheduling.}
%\end{abstract}
%
%
%

\section{Contributions}
As a team, we worked very well and organised throughout the semester. My project partner offered valuable industry experience and established development practices, while I contributed organisational and planning skills that helped us keep the project aligned with our broader goals. I regularly circled back to open questions and ensured the same understanding of objectives and implementation details.  

We worked on most of the contributions in this project together in collaboration sessions, such as the adjustment of the Python pipeline for track malfunctions, debugging and the basic version of the secondary encoding. I contributed the malfunction log and created the testing environments. Before the main experiments, I ran some preliminary tests. 
Due to my Linguistics background, I could provide experience in academic writing, reviewing literature and typesetting documents with \LaTeX. I wrote and refined substantial parts of the final paper and the presentation slides.


\section{Reflection}
I am very pleased with the process, experiment results and outcomes of this project. It allowed me to improve in mutliple essential skills, such as pipeline building and adjustment, Python and Clingo programming and use of Git.

The project benefited from the complementary skills we brought as a team. One personal challenge was working across different areas of experience and confidence. My teammate’s background in software development provided valuable expertise in programming and helped advance the technical implementation, while my own background contributed organisational, academic-writing, and communication skills. This required me to become more confident in contributing to areas that were initially less familiar to me. Through the project, particularly while preparing the final paper and presentation, I gained a clearer appreciation of the value of my own contributions and developed greater confidence in my abilities.


\section{Recommendation}
I would recommend future students to spend some time discussing how to work together on a project and which skills they can bring in. In my experience, this can be a tricky subject to bring up because it requires openness and vulnerability to discuss personal preferences and dislikes. Nevertheless, I would encourage every student to take this course and pick a project partner with a different skill set than their own because it offers an excellent opportunity to advance your own skills.


\end{document}
